\documentclass[11pt]{article}

\usepackage[margin=1in]{geometry}
\usepackage{amsmath,amssymb,amsthm,mathtools}
\usepackage{enumitem}
\usepackage{xcolor}
\usepackage[colorlinks=true,linkcolor=black,urlcolor=blue]{hyperref}
\usepackage{fancyhdr}

\pagestyle{fancy}
\setlength{\headheight}{14pt}
\fancyhf{}
\lhead{CMPUT 654: Mathematical Foundations of Modern AI}
\rhead{Homework 0}
\cfoot{\thepage}

\newcommand{\R}{\mathbb{R}}
\newcommand{\E}{\mathbb{E}}
\newcommand{\Pp}{\mathbb{P}}
\newcommand{\norm}[1]{\left\lVert #1\right\rVert}
\DeclareMathOperator*{\argmin}{arg\,min}
\DeclareMathOperator*{\argmax}{arg\,max}

\newtheoremstyle{courseproblem}
  {1.1em}{1.1em}{\normalfont}{}{\bfseries}{.}{0.6em}{}
\theoremstyle{courseproblem}
\newtheorem{problem}{Problem}

\begin{document}

\begin{center}
{\Large\bfseries Homework 0: Mathematical Readiness}\\[0.35em]
{\large CMPUT 654, Fall 2026}
\end{center}

\paragraph{Purpose.}
This assignment is not graded. It samples the proof writing, probability,
linear algebra, calculus, and abstraction used in the course. Attempt it during
the first week without searching for solutions. If substantial parts remain
inaccessible after reviewing the relevant prerequisites, speak with the
instructor promptly and make an honest decision about whether the course is a
good fit.

\paragraph{Submission.}
No submission is required unless the instructor announces otherwise. Students
who would like diagnostic feedback may submit a typed PDF by the end of Week 1.
Use the course homework template, acknowledge every source and collaborator,
and write every argument in your own words.

\paragraph{Total.} 100 diagnostic points. The points indicate intended effort,
not a course grade.

\begin{problem}[Log loss identifies a conditional probability --- 25 points]
Let $Y\sim\operatorname{Bernoulli}(\eta)$, where $\eta\in(0,1)$. For a
prediction $q\in(0,1)$, define the log loss
\[
\ell(q;Y)=-Y\log q-(1-Y)\log(1-q)
\]
and its population risk $L_\eta(q)=\E[\ell(q;Y)]$.
\begin{enumerate}[label=(\alph*)]
  \item Compute $L_\eta'(q)$ and $L_\eta''(q)$, and prove that $q=\eta$ is
  its unique minimizer. \hfill [7]
  \item Show that
  \[
  L_\eta(q)-L_\eta(\eta)
  =\eta\log\frac{\eta}{q}+(1-\eta)\log\frac{1-\eta}{1-q}.
  \]
  Identify the expression on the right. \hfill [6]
  \item You may use the Bernoulli form of Pinsker's inequality,
  $D_{\mathrm{KL}}(\operatorname{Ber}(\eta)\|\operatorname{Ber}(q))
  \ge 2(\eta-q)^2$. Deduce a bound on $|q-\eta|$ from an upper bound
  $L_\eta(q)-L_\eta(\eta)\le\varepsilon$. \hfill [5]
  \item A classifier predicts $1$ when $q\ge 1/2$ and $0$ otherwise. Find
  the Bayes-optimal classifier under zero--one loss and explain why a small
  probability-estimation error need not guarantee the correct classification
  when $\eta$ is close to $1/2$. \hfill [7]
\end{enumerate}
\end{problem}

\begin{problem}[Uniform estimation and selection --- 20 points]
For every $j\in[m]$, let $Z_{j,1},\ldots,Z_{j,n}$ be independent random
variables taking values in $[0,1]$, with common mean $\mu_j$ for fixed $j$.
Let $\widehat\mu_j=n^{-1}\sum_{i=1}^n Z_{j,i}$. You may use Hoeffding's
inequality
\[
\Pp\!\left(|\widehat\mu_j-\mu_j|\ge\varepsilon\right)
\le 2e^{-2n\varepsilon^2}.
\]
\begin{enumerate}[label=(\alph*)]
  \item Prove that, with probability at least $1-\delta$,
  $\max_{j\in[m]}|\widehat\mu_j-\mu_j|\le\varepsilon$, provided
  \[
  n\ge \frac{\log(2m/\delta)}{2\varepsilon^2}.
  \]
  State exactly where independence is used and where it is not used. \hfill [10]
  \item Let $\widehat j\in\argmax_j\widehat\mu_j$ and
  $j^\star\in\argmax_j\mu_j$. On the event above, prove
  $\mu_{j^\star}-\mu_{\widehat j}\le 2\varepsilon$. \hfill [10]
\end{enumerate}
\end{problem}

\begin{problem}[Which interpolator does gradient descent select? --- 30 points]
Let $A\in\R^{n\times d}$ have full row rank, with $n<d$, and let
$b\in\R^n$. Consider
\[
f(w)=\frac12\norm{Aw-b}_2^2
\]
and gradient descent
\[
w_{t+1}=w_t-\gamma A^\top(Aw_t-b),\qquad w_0=0.
\]
\begin{enumerate}[label=(\alph*)]
  \item Show that
  $w^\dagger=A^\top(AA^\top)^{-1}b$ interpolates the data and is the unique
  minimum-Euclidean-norm vector satisfying $Aw=b$. \hfill [10]
  \item Prove that every iterate $w_t$ lies in the row space of $A$. \hfill [5]
  \item Let $r_t=Aw_t-b$. Show
  $r_{t+1}=(I-\gamma AA^\top)r_t$. Prove that $r_t\to0$ whenever
  $0<\gamma<2/\norm{A}_2^2$. \hfill [10]
  \item Conclude that $w_t\to w^\dagger$. What can change if $w_0$ has a
  nonzero component in the null space of $A$? \hfill [5]
\end{enumerate}
\end{problem}

\begin{problem}[A small query-complexity lower bound --- 25 points]
An unknown index $i^\star\in[k]$ is selected. On query $j\in[k]$, a learner
observes $1$ if $j=i^\star$ and $0$ otherwise. The learner must eventually
identify $i^\star$ with probability one. It may choose queries adaptively and
may randomize.
\begin{enumerate}[label=(\alph*)]
  \item Give a deterministic learner making at most $k-1$ queries on every
  instance, and prove it is correct. \hfill [6]
  \item Prove that every deterministic learner requires at least $k-1$
  queries on some instance. \hfill [6]
  \item Consider querying indices in a uniformly random order and stopping
  after finding $i^\star$, or after eliminating $k-1$ indices. Show that for
  every $i^\star$, its expected number of queries is
  \[
  \frac{k+1}{2}-\frac1k.
  \]
  \hfill [6]
  \item Use Yao's minimax principle with a uniform random $i^\star$ to prove
  that every randomized learner has worst-case expected query complexity at
  least $(k+1)/2-1/k$. State the version of Yao's principle you use. \hfill [7]
\end{enumerate}
\end{problem}

\end{document}
